\begin{lemma}[Pairing graph]
\label{lem:pairing-graph}
Let
\[
  \mathsf P(x,y,z)\defeq(x+y)\times S(x+y)+(y+y)=z+z.
\]
Then \(\mathsf P\) is a quantifier-free \(\Lang_A\)-formula; for all
\(a,b,c\in\N\), \(\mathcal N\entails\mathsf P(\bar a,\bar b,\bar c)\) iff
\(c=\operatorname{pair}_{\N}(a,b)\); and both
\(a\leq\operatorname{pair}_{\N}(a,b)\) and
\(b\leq\operatorname{pair}_{\N}(a,b)\) hold for all \(a,b\in\N\).
\textit{Status: complete proof.}
\end{lemma}

\begin{proof}
The display is a single atomic equation, so \(\mathsf P\) is
quantifier-free.  In \(\mathcal N\) it evaluates to
\((a+b)(a+b+1)+2b=2c\), which holds iff
\(c=(a+b)(a+b+1)/2+b=\operatorname{pair}_{\N}(a,b)\); the doubled form
avoids division, which \(\Lang_A\) does not provide.  For the monotonicity
bounds: \(b\leq(a+b)(a+b+1)/2+b\) is immediate, and
\(a\leq a+b\leq(a+b)(a+b+1)/2\leq\operatorname{pair}_{\N}(a,b)\), where the
middle inequality holds because \(a+b=0\) forces \(a=0\), while
\(a+b\geq1\) makes the factor \((a+b+1)/2\) at least \(1\).
\end{proof}

In the next display, \(d\) abbreviates the term \(S((Si)\times m)\), and
\(u\) is a fresh variable with \(u\notin\operatorname{Var}(d)\).  The final
conjunct \(\exists u<d\,(v+Su=d)\) is the \emph{marked-bounded} form of the
comparison \(v<d\): it is \(\mathcal N\)-equivalent to the plain
abbreviation of \(v<d\), because any witness of \(v+Su=d\) satisfies
\(Su\leq d\), hence \(u<d\).  The marked form is required because the
hierarchy definition above counts only marked bounded quantifiers as
bounded, while the bare comparison \(v<d\) expands to an unbounded
existential.

\begin{lemma}[G\"odel beta function]
\label{lem:goedel-beta}
Define
\[
  \mathsf B(w,i,v)\defeq\exists t<Sw\,\exists m<Sw\,\exists q<St\,
  \bigl(\mathsf P(t,m,w)\land t=d\times q+v\land
  \exists u<d\,(v+Su=d)\bigr).
\]
Then \(\mathsf B\) is a bounded arithmetic formula, and:
\begin{enumerate}[label=(\alph*)]
\item for all \(w,i\in\N\) there is exactly one \(v\in\N\) with
      \(\mathcal N\entails\mathsf B(\bar w,\bar i,\bar v)\), and this
      \(v\) satisfies \(v\leq w\); write \(\beta(w,i)\) for it; and
\item for every finite sequence \(a_0,\ldots,a_{k-1}\) of natural numbers
      there is \(w\in\N\) with \(\beta(w,i)=a_i\) for every \(i<k\).
\end{enumerate}
\textit{Status: complete proof.}
\end{lemma}

\begin{proof}
Boundedness: every quantifier in the display is marked bounded, the final
conjunct is the marked-bounded form of \(v<d\) as recorded above, and the
remaining conjuncts are atomic.

(a) Since \(\operatorname{pair}_{\N}\) is a bijection, exactly one pair
\((t,m)\) satisfies \(\operatorname{pair}_{\N}(t,m)=w\), and
\cref{lem:pairing-graph} gives \(t\leq w\) and \(m\leq w\), so the
displayed bounds \(t<Sw\) and \(m<Sw\) do not exclude this pair, while
\(\mathsf P(t,m,w)\) excludes every other.  The term \(d\) evaluates to
\(1+(i+1)m\geq1\), so Euclidean division of \(t\) by that value determines
\(q\) and \(v\) uniquely via \(t=(1+(i+1)m)\,q+v\) with \(v<1+(i+1)m\);
moreover \(q\leq t<St\) and \(v\leq t\leq w\), so the bound on \(q\) is
respected and \(v\leq w\).

(b) Given \(a_0,\ldots,a_{k-1}\), put
\(N\defeq\max(k,a_0,\ldots,a_{k-1})+1\) and \(m\defeq N!\).  The moduli
\(1+(i+1)m\) for \(i<k\) are pairwise coprime: a prime \(p\) dividing both
\(1+(i+1)m\) and \(1+(j+1)m\) with \(i<j<k\) divides their difference
\((j-i)m\) but not \(m\), hence divides \(j-i\); yet \(0<j-i<k\leq N\)
gives \(p\leq N\), so \(p\) divides \(N!=m\), a contradiction.  Each
\(a_i<N\leq m<1+(i+1)m\).  By the Chinese remainder theorem for pairwise
coprime moduli there is \(t\) with \(t\equiv a_i\pmod{1+(i+1)m}\) for
every \(i<k\), and since \(a_i<1+(i+1)m\), the remainder of \(t\) on
division by \(1+(i+1)m\) is exactly \(a_i\).  Put
\(w\defeq\operatorname{pair}_{\N}(t,m)\); then \(\beta(w,i)=a_i\) for
every \(i<k\).
\end{proof}

G\"odel's original beta-function argument \cite{Godel1931} and the
treatment in \cite[ch.~16]{BoolosBurgessJeffrey2007} are parallel reading
for \cref{lem:goedel-beta}.

\begin{lemma}[Pairing contraction]
\label{lem:pair-contraction}
Let \(\theta\) be an \(\Lang_A\)-formula and let \(z\) be a variable not
free in \(\theta\) and distinct from \(x,y\).  Then \(\mathcal N\)
satisfies, under every variable assignment,
\[
  \exists x\,\exists y\,\theta
  \;\leftrightarrow\;
  \exists z\,\exists x<Sz\,\exists y<Sz\,(\mathsf P(x,y,z)\land\theta),
\]
and dually
\[
  \forall x\,\forall y\,\theta
  \;\leftrightarrow\;
  \forall z\,\forall x<Sz\,\forall y<Sz\,(\mathsf P(x,y,z)\to\theta).
\]
\textit{Status: complete proof.}
\end{lemma}

\begin{proof}
Fix an assignment.  If \(a,b\in\N\) satisfy \(\theta\) at \(x,y\), put
\(c\defeq\operatorname{pair}_{\N}(a,b)\); then \(\mathsf P(a,b,c)\) holds
in \(\mathcal N\) by \cref{lem:pairing-graph}, and its monotonicity bounds
give \(a,b\leq c<c+1\), so \(a\) and \(b\) witness the two marked bounded
quantifiers under \(z\mapsto c\).  Conversely, witnesses of the right side
satisfy \(\theta\) in particular.  Since \(z\) is not free in \(\theta\)
and differs from \(x,y\), no capture occurs and \(\theta\) keeps its
value.  The dual equivalence is the first applied to \(\neg\theta\),
negating both sides and moving negations inward.
\end{proof}

\begin{lemma}[Bounded-quantifier exchange]
\label{lem:bounded-collection}
Let \(\theta\) be an \(\Lang_A\)-formula, let \(t\) be an
\(\Lang_A\)-term with \(y,w\notin\operatorname{Var}(t)\), and let \(w\) be
fresh.  Then \(\mathcal N\) satisfies, under every variable assignment,
\[
  \forall y<t\,\exists x\,\theta
  \;\leftrightarrow\;
  \exists w\,\forall y<t\,\exists x<w\,\theta,
  \qquad
  \exists y<t\,\forall x\,\theta
  \;\leftrightarrow\;
  \forall w\,\exists y<t\,\forall x<w\,\theta.
\]
Both are equivalences of satisfaction in the standard structure, used
below for the boldface definability classes only; they are not
derivability claims in any formal theory of arithmetic.
\textit{Status: complete proof.}
\end{lemma}

\begin{proof}
Fix an assignment and consider the first equivalence.  Right to left is
immediate: a witness below \(w\) is in particular a witness.  Left to
right: the assignment gives \(t\) a value \(n\), and for each of the
finitely many \(b<n\) there is by hypothesis a witness \(x_b\in\N\) for
\(\theta\) at \(y\mapsto b\); let \(c\) be \(0\) when \(n=0\) and
\(\max_{b<n}x_b+1\) otherwise.  Then \(w\mapsto c\) makes
\(\forall y<t\,\exists x<w\,\theta\) true.  Freshness of \(w\) and
\(y,w\notin\operatorname{Var}(t)\) prevent capture and keep the value of
\(t\) fixed.  The second equivalence is the first applied to
\(\neg\theta\), negating both sides and moving negations inward.
\end{proof}

Call two \(\Lang_A\)-formulas \emph{\(\mathcal N\)-equivalent} iff exactly
the same variable assignments satisfy them in \(\mathcal N\).  Call a
formula \emph{prefix-form} \(\Sigma^0_n\) iff it consists of \(n\)
alternating blocks of unbounded quantifiers, the first block existential,
applied to a bounded arithmetic formula in negation normal form, and
\emph{prefix-form} \(\Pi^0_n\) dually, with the first block universal.
Every prefix-form formula lies in the corresponding syntactic class, since
the generating clauses allow each unbounded quantifier to be prepended at
the appropriate level.  The cited hierarchy sources define their classes
by prefix forms; the next theorem identifies the two presentations.

\begin{theorem}[Prefix normal form]
\label{thm:prefix-normal-form}
For every \(n\geq0\):
\begin{enumerate}[label=(\alph*)]
\item the relations defined in \(\mathcal N\) by prefix-form
      \(\Sigma^0_n\) formulas are closed, up to \(\mathcal N\)-equivalence
      of the defining formulas, under \(\land\), \(\lor\), and both marked
      bounded quantifiers, and for \(n\geq1\) also under unbounded
      existential quantification; dually for \(\Pi^0_n\) with unbounded
      universal quantification;
\item every \(\Sigma^0_n\) formula is \(\mathcal N\)-equivalent to a
      prefix-form \(\Sigma^0_n\) formula, and every \(\Pi^0_n\) formula to
      a prefix-form \(\Pi^0_n\) formula; and
\item the boldface classes \(\boldsymbol{\Sigma}^0_n(k)\),
      \(\boldsymbol{\Pi}^0_n(k)\), and \(\boldsymbol{\Delta}^0_n(k)\)
      coincide with the classes of relations defined in \(\mathcal N\) by
      prefix-form formulas of the corresponding kinds.
\end{enumerate}
\textit{Status: complete proof.}
\end{theorem}

\begin{proof}
Two preliminary observations.  \emph{Padding}: for \(m<n\), every
prefix-form formula of level \(m\) is \(\mathcal N\)-equivalent to a
prefix-form \(\Sigma^0_n\) formula and to a prefix-form \(\Pi^0_n\)
formula, obtained by prepending and interleaving quantifiers on fresh
variables; each added quantifier is vacuous because the domain \(\N\) is
nonempty.  \emph{Normal-form bookkeeping}: an implication
\(\mathsf P\to\theta\) introduced by the dual contraction of
\cref{lem:pair-contraction} is written \(\neg\mathsf P\lor\theta\), which
is in negation normal form because \(\mathsf P\) is atomic.

(a) Induction on \(n\), both classes at once; the \(\Pi^0_n\) case is dual
throughout, exchanging \(\exists\) with \(\forall\), \(\land\) with
\(\lor\) where the argument negates, and
\cref{lem:pair-contraction,lem:bounded-collection} with their dual halves.

Level \(n=0\): a prefix-form \(\Sigma^0_0\) or \(\Pi^0_0\) formula is a
bounded formula in negation normal form, and \(\land\), \(\lor\), and
marked bounded quantification produce bounded formulas in negation normal
form syntactically; no unbounded closure is claimed at this level.

Level \(n\geq1\): a prefix-form \(\Sigma^0_n\) formula is
\(\mathcal N\)-equivalent, by iterating the contraction of
\cref{lem:pair-contraction} inside its first block, to one of the form
\(\exists x\,B\) with \(B\) prefix-form \(\Pi^0_{n-1}\); the residual
bounded quantifiers and the conjuncts \(\mathsf P\) introduced by each
contraction are absorbed into \(B\) by the induction hypothesis for level
\(n-1\).  It therefore suffices to treat formulas \(\exists x\,B\) and
\(\exists x'\,B'\) with \(B,B'\) prefix-form \(\Pi^0_{n-1}\).
\begin{enumerate}[label=(\roman*)]
\item \(\land\), \(\lor\): rename the bound variables apart, so that
      \(x\notin\FV(B')\) and \(x'\notin\FV(B)\), and pull the existentials
      out: \(\exists x\,B\land\exists x'\,B'\) is
      \(\mathcal N\)-equivalent to \(\exists x\,\exists x'(B\land B')\),
      and \(\exists x\,B\lor\exists x'\,B'\) to
      \(\exists x\,\exists x'(B\lor B')\), the latter because the
      quantifier on the disjunct that did not supply the witness is
      vacuous over the nonempty domain.  Contract the two existentials by
      \cref{lem:pair-contraction}; the resulting matrix applies two marked
      bounded quantifiers and a conjunction with the quantifier-free
      \(\mathsf P\) to \(B\land B'\), respectively \(B\lor B'\), each
      \(\mathcal N\)-equivalent to a prefix-form \(\Pi^0_{n-1}\) formula
      by the induction hypothesis; prepending \(\exists z\) gives
      prefix-form \(\Sigma^0_n\).
\item Marked bounded \(\exists\): with \(x\) renamed so that \(x\neq y\)
      and \(x\notin\operatorname{Var}(t)\),
      \(\exists y<t\,\exists x\,B\) is \(\mathcal N\)-equivalent to
      \(\exists x\,\exists y<t\,B\) by pure logic; the induction
      hypothesis absorbs \(\exists y<t\) into \(B\), and prepending
      \(\exists x\) gives prefix-form \(\Sigma^0_n\).
\item Marked bounded \(\forall\): by \cref{lem:bounded-collection},
      \(\forall y<t\,\exists x\,B\) is \(\mathcal N\)-equivalent to
      \(\exists w\,\forall y<t\,\exists x<w\,B\); the induction hypothesis
      absorbs the two marked bounded quantifiers into \(B\), and
      prepending \(\exists w\) gives prefix-form \(\Sigma^0_n\).
\item Unbounded \(\exists\): contract
      \(\exists x'\,\exists x\,B\) by \cref{lem:pair-contraction} and
      absorb as in (i).
\end{enumerate}

(b) By the hierarchy definition, an arbitrary formula is \(\Sigma^0_n\)
or \(\Pi^0_n\) exactly when its fixed negation-normal-form conversion
lies in the corresponding generated class, and each conversion step ---
implication elimination, the De Morgan and quantifier-duality rewrites,
the bounded dualities, and capture-avoiding renaming --- preserves
satisfaction under every assignment, so every formula is
\(\mathcal N\)-equivalent to its conversion.  It therefore suffices to
prove the claim for the members of the generated classes, by induction on
\(n\) and, within each level, by induction on the generating clauses of
the syntactic hierarchy definition.  At level \(0\) the generated classes
consist of the bounded formulas in negation normal form, which are
exactly the prefix-form formulas of level \(0\).  At level \(n+1\), the
base formulas are the
\(\Sigma^0_n\) and \(\Pi^0_n\) formulas, which are
\(\mathcal N\)-equivalent to prefix-form formulas of level \(n\) by the
induction hypothesis and hence, by padding, to prefix-form
\(\Sigma^0_{n+1}\) formulas; each closure clause --- \(\land\), \(\lor\),
either marked bounded quantifier, unbounded \(\exists\) --- preserves
\(\mathcal N\)-equivalence to prefix form by (a) at level \(n+1\).  The
\(\Pi^0_{n+1}\) half is dual.

(c) By (b), a relation definable by a \(\Sigma^0_n\) formula is definable
by a prefix-form \(\Sigma^0_n\) formula, and conversely every prefix-form
\(\Sigma^0_n\) formula is a \(\Sigma^0_n\) formula; likewise for
\(\Pi^0_n\), hence for the intersections \(\boldsymbol{\Delta}^0_n(k)\).
\end{proof}

